\section{Conclusion}

Real-world agent work makes timeline views more important and more dangerous.
They are important because humans and agents need durable surfaces for recovery,
audit, and cooperation. They are dangerous when they hide their perspective and
present mixed-source facts as an absolute global order.

Observer-declared timelines offer a practical alternative. A product can keep
facts and causality load-bearing while admitting that every useful view is
projected from a declared perspective. This paper adds a further constraint:
the observer is not timeless. A stable logical observer may have many
incarnations, and continuity between them must be established through accepted
cuts, policy/runtime identity, recovery evidence, and explicit degradation.

Kungfu now supplies preliminary implementation evidence for both sides of this
rule. Episodes provide bounded causal authority and safe-capability evidence.
Exact-cut runtime readiness, generation and coordinator-epoch fencing, Peer
rebootstrap, explicit recovery states, and immutable runtime selection show how
an observer incarnation can fail or change without silently acquiring old
authority. Historical scale evidence and newer source-bound single-host
qualification remain separate, preserving the provenance of each claim.

The results do not yet validate the complete observer-relative multi-machine
view. That next step requires first-class observer policy, accepted source
ranges, reproducible cross-source ordering, exported incarnation and continuity
declarations, projection-level causal fsck, and experiments in which humans and
agents recover from the same declared perspective. Keeping that boundary
explicit is part of the proposal itself: trust in a timeline, including trust in
its continuity, must start from the facts actually available.
